\section{DelegNet}
\label{sec:method}

\subsection{Problem formulation}
A delegation event $e=(u,s,c,t)$ records that authority $u$ delegated to contract $c$ through a transaction submitted by sponsor $s$ at time $t$. Let $\mathcal{H}_{\tau}$ denote all chain events observable up to a decision time $\tau$. A detector is a function $f(e,\mathcal{H}_{\tau}) \in [0,1]$ estimating the probability that $c$ is malicious. We consider three decision times:
\begin{itemize}
\item \emph{pre-signature}, $\tau=t^-$: the authorization is not yet on chain and the sponsor is unknown, which is the situation of a wallet asked to sign;
\item \emph{inclusion}, $\tau=t$: the authorization and its transaction are visible;
\item \emph{monitoring}, $\tau=t+\delta$, $\delta\in\{12,60,300,1800,7200\}$\,s.
\end{itemize}
The defender's training data contain labelled contracts from code families $\mathcal{F}_{\mathrm{train}}$, and evaluation contracts come from disjoint families $\mathcal{F}_{\mathrm{test}}$, $\mathcal{F}_{\mathrm{train}}\cap\mathcal{F}_{\mathrm{test}}=\emptyset$.

\subsection{Design rationale}
Two observations drive the design. First, the semantics that a theft requires---forwarding the account's value to attacker-controlled destinations, or granting an attacker-controlled caller arbitrary execution---must be present in the delegate's code, whereas the behavioural context of a delegation (who submits it, in what batch, from which funding source, how long before the drain) is chosen freely and cheaply by the attacker. A detector whose decision rests on code semantics therefore resists the cheapest evasions. Second, attackers mass-produce their contracts by copying, and the copies differ by exactly the edits that defeat syntactic matching: new hard-coded addresses, junk instructions, equivalent instruction choices and appended data. A representation that is invariant to those edits, but preserves what the code does, turns copy-paste into a liability for the attacker rather than an evasion.

\dn{} is organised as a pattern-recognition module with three stages: (A)~a canonical code representation that removes attacker-controlled syntactic variation, (B)~interpretable theft idioms computed on that representation, and (C)~a family-exclusive memory of known scam and benign code patterns that produces the detection score and, through its nearest match, a case-based explanation.

\subsection{Canonical code representation}
\label{sec:canon}
Delegate contracts execute on the Ethereum Virtual Machine (EVM), a stack machine whose bytecode interleaves single-byte opcodes and immediate operands (\texttt{PUSH1}--\texttt{PUSH32} carry 1--32 operand bytes). Control flow uses \texttt{JUMP}/\texttt{JUMPI} to absolute byte offsets that must hold a \texttt{JUMPDEST} instruction. Our canonicalisation maps bytecode $b$ to a token sequence $\kappa(b)$ in five steps.

\paragraph{Step 1: disassembly and code-address identification} The bytecode is decoded into instructions $(o_i,v_i,\pi_i)$ with opcode $o_i$, immediate $v_i$ and byte offset $\pi_i$. Undefined bytes are kept verbatim as \texttt{INVALID} with their raw value, so that decoding and re-encoding are lossless. We identify immediates that encode code addresses: a push whose value equals the offset of a \texttt{JUMPDEST} and is immediately followed by a jump (a static jump), or a push of at least two bytes whose value equals a \texttt{JUMPDEST} offset (a return address pushed for an internal call, the pattern Solidity emits for function calls). One-byte pushes are tagged only if their value is not a common memory or calldata constant (multiples of 32 and the calldata argument offsets $4+32k$), because those values coincide with low code offsets in small contracts.

\paragraph{Step 2: reachability} Starting from offset $0$, we compute the set of reachable instructions by a worklist closure. Execution falls through until a terminator (\texttt{STOP}, \texttt{RETURN}, \texttt{REVERT}, \texttt{INVALID}, \texttt{SELFDESTRUCT}, \texttt{JUMP}), and every \texttt{JUMPDEST} whose offset is pushed by a reachable instruction becomes a new root. Unreachable instructions---compiler metadata, appended data and attacker-appended padding---are dropped.

\paragraph{Step 3: equivalence normalisation} Instruction sequences with identical effect on the stack and state are rewritten to one representative: \texttt{ADDRESS BALANCE}$\,\to\,$\texttt{SELFBALANCE}; \texttt{XOR ISZERO}$\,\to\,$\texttt{EQ} (since $a=b \Leftrightarrow (a \oplus b)=0$); \texttt{PUSH1 0x00}$\,\to\,$\texttt{PUSH0}.

\paragraph{Step 4: junk elimination} A peephole pass removes, to a fixed point, sequences whose net effect on the stack is empty and which have no side effect: a push followed by \texttt{POP}; an environment read (\texttt{GAS}, \texttt{CALLVALUE}, \texttt{PC}, \texttt{MSIZE}, \texttt{ADDRESS}, \texttt{CALLER}, \texttt{CALLDATASIZE}, \texttt{RETURNDATASIZE}) followed by \texttt{POP}; \texttt{PUSH0 ISZERO POP}; and \texttt{PUSH}$\,$\texttt{PUSH}$\,$\texttt{SWAP1}$\,$\texttt{POP}$\,$\texttt{POP}. \texttt{JUMPDEST} markers are removed entirely, because their positions are precisely what instruction insertion shifts, while the control-flow information they carry is preserved by Step~5.

\paragraph{Step 5: semantic tagging} Pushed code addresses and all two- and three-byte pushes become a single \texttt{TAG} token, so that re-laid-out jump targets compare equal. Four-byte pushes equal to well-known function selectors (ERC-20 \texttt{transfer}, \texttt{transferFrom}, \texttt{approve}, \texttt{setApprovalForAll}, \texttt{permit}, Permit2, ERC-1271 \texttt{isValidSignature}, ERC-4337 \texttt{validateUserOp}, \texttt{execute} variants, proxy upgrades, token-receiver hooks) become named tokens. Twenty-byte pushes become \texttt{ADDR}, so that copies with different hard-coded owners compare equal. Small constants in $[1,10]$---including precompile addresses such as \texttt{ecrecover}---are tagged by push width. All other opcodes are kept.

\begin{algorithm}[t]
\caption{Canonicalisation $\kappa(b)$ of EVM bytecode}
\label{alg:canon}
\begin{algorithmic}[1]
\REQUIRE bytecode $b$
\STATE $I \leftarrow$ decode$(b)$ \COMMENT{lossless; undefined bytes kept}
\STATE $D \leftarrow \{\pi_i : o_i=\texttt{JUMPDEST}\}$
\STATE mark $i$ as code-address push if $v_i\in D$ and ($o_{i+1}\in\{\texttt{JUMP},\texttt{JUMPI}\}$ or width $\ge 2$ or $v_i\notin C_{\text{common}}$)
\STATE $R \leftarrow \emptyset$; $Q \leftarrow \{0\}$ \COMMENT{reachability}
\WHILE{$Q\neq\emptyset$}
\STATE pop $k$; \textbf{while} $k\notin R$: add $k$ to $R$; push target of any reachable code-address push onto $Q$; \textbf{break} at terminator; $k\leftarrow k+1$
\ENDWHILE
\STATE $I \leftarrow [\,i\in I : i\in R\,]$
\STATE rewrite equivalences: \texttt{ADDRESS BALANCE}$\to$\texttt{SELFBALANCE}, \texttt{XOR ISZERO}$\to$\texttt{EQ}, \texttt{PUSH1 0}$\to$\texttt{PUSH0}
\STATE drop \texttt{JUMPDEST}; remove stack-neutral suffixes ending in \texttt{POP} to a fixed point
\STATE map: code-address and 2--3-byte pushes $\to$ \texttt{TAG}; 20-byte pushes $\to$ \texttt{ADDR}; known selectors $\to$ names; small constants $\to$ width tags
\RETURN token sequence
\end{algorithmic}
\end{algorithm}

\begin{table}[t]
\centering
\caption{Worked example: the forwarding core of a sweeper, a disguised copy, and their common canonical form. Junk (italic) and equivalent substitutions (underlined) are removed or normalised.}
\label{tab:canon-example}
\footnotesize
\setlength{\tabcolsep}{2pt}
\begin{tabular}{p{0.3\columnwidth}p{0.36\columnwidth}p{0.26\columnwidth}}
\toprule
Original & Disguised copy & Canonical \\
\midrule
\texttt{PUSH0} & \texttt{PUSH1 0x00} & \texttt{PUSH0} \\
\texttt{PUSH0 PUSH0 PUSH0} & \emph{\texttt{GAS POP}} \texttt{PUSH0 PUSH0 PUSH0} & \texttt{PUSH0 PUSH0 PUSH0} \\
\texttt{SELFBALANCE} & \underline{\texttt{ADDRESS BALANCE}} & \texttt{SELFBALANCE} \\
\texttt{PUSH20 0xA1\ldots} & \texttt{PUSH20 0x7C\ldots} (new owner) & \texttt{ADDR} \\
\texttt{GAS CALL POP} & \emph{\texttt{PUSH2 0x3f1 PUSH1 9 SWAP1 POP POP}} \texttt{GAS CALL POP} & \texttt{GAS CALL POP} \\
\texttt{STOP} & \texttt{STOP} \emph{\texttt{INVALID 0x5b\ldots}} (trailer) & \texttt{STOP} \\
\bottomrule
\end{tabular}
\end{table}

Algorithm~\ref{alg:canon} summarises the procedure and Table~\ref{tab:canon-example} illustrates it on the forwarding core shared by most sweepers in our data. Raw four-gram similarity between the two left columns is low---every inserted instruction breaks several four-grams---whereas their canonical sequences are identical. The idiom ``\texttt{SELFBALANCE} feeding a \texttt{CALL} to a hard-coded address'' survives intact.

\paragraph{Invariance} Let $\mathcal{M}$ be the set of mutations that (i)~insert a stack-neutral sequence from the set $\mathcal{N}$ eliminated in Step~4 at any instruction boundary other than between a static push and its jump, (ii)~apply a rewrite from Step~3 in either direction, (iii)~append an unreachable trailer, and (iv)~re-encode code-address immediates to the new offsets, possibly widening the push. For any $b$ and $m\in\mathcal{M}^{*}$ that preserves the reachable control-flow graph, $\kappa(m(b))=\kappa(b)$. The argument is constructive. Trailers are unreachable and removed by Step~2. Inserted $\mathcal{N}$-sequences are reachable but eliminated by Step~4, whose rules are confluent because each rule removes a contiguous suffix ending in \texttt{POP} and introduces no new instructions. Rewrites from Step~3 map to the same representative. Re-encoded code addresses become \texttt{TAG} irrespective of value and width, and \texttt{JUMPDEST} tokens are removed, so offset shifts are invisible. Equality fails only when a mutation creates a sequence outside $\mathcal{N}$ or alters the reachable control flow, which by definition changes the program. Empirically, over $60$ malicious contracts and six mutation strengths, the canonical four-gram similarity between a contract and its mutated copy is $0.999$ on average, against $0.192$--$0.766$ for raw opcodes (Section~\ref{sec:res-disguise}).

\paragraph{Mutation engine} To evaluate and train against disguise we implemented a semantics-preserving mutator over disassembled code. It inserts stack-neutral sequences with probability $\min(0.5,0.06\,\sigma)$ per instruction, substitutes equivalent forms with probability $\min(0.9,0.25\,\sigma)$, appends an unreachable trailer of up to $24\sigma$ bytes, and re-links all identified code addresses, iteratively widening pushes until every target fits. Mutation strength $\sigma$ controls all three. Each mutant is verified: its reachable instruction count must not decrease, otherwise it is regenerated (after five failures the mutator falls back to substitution only). This guarantees that evaluation uses only mutants whose control flow survives.

\subsection{Code idioms}
From $\kappa(b)$ we compute $21$ interpretable indicators of theft-relevant semantics:
\begin{itemize}
\item value forwarding: a \texttt{CALL} preceded within twelve instructions by \texttt{SELFBALANCE}, and specifically to a hard-coded \texttt{ADDR};
\item caller authorisation: \texttt{CALLER} compared with a hard-coded address (owner gate) or with \texttt{ADDRESS} (self-call gate);
\item \texttt{ORIGIN} use;
\item signature verification: a \texttt{STATICCALL} to the \texttt{ecrecover} precompile;
\item \texttt{DELEGATECALL} in general, and to a calldata- or storage-derived target;
\item \texttt{SELFDESTRUCT};
\item token-transfer, transfer-from, approval, permit, upgrade and account-abstraction selectors;
\item counts of dispatched selectors, hard-coded addresses, calls and storage writes;
\item canonical length, and whether the fallback region performs a call.
\end{itemize}
Because they are computed on $\kappa(b)$, the idioms inherit its invariance. Over all evaluated contracts, $34$--$36$ of $36$ contracts retained identical idiom vectors under mutation at every strength during development.

\subsection{Code-pattern memory and detection score}
\label{sec:memory}
Let $G(b)$ be the set of four-grams of $\kappa(b)$, and $J(A,B)=|A\cap B|/|A\cup B|$ the Jaccard similarity. From the training delegations we build a memory $\mathcal{M}=\{(F_j,G_j,\bar y_j)\}$ containing, for each training contract $j$, its code family $F_j$, four-gram set $G_j$ and mean label $\bar y_j$. For a query contract $q$ with family $F_q$ we compute
\begin{align}
\sigma_{\mathrm{mal}}(q)&=\max_{j:\,\bar y_j\ge 0.5,\;F_j\ne F_q} J\big(G(q),G_j\big),\\
\sigma_{\mathrm{ben}}(q)&=\max_{j:\,\bar y_j< 0.5,\;F_j\ne F_q} J\big(G(q),G_j\big),
\end{align}
and the detection score $f(q)=\tfrac{1}{2}\big(1+\sigma_{\mathrm{mal}}(q)-\sigma_{\mathrm{ben}}(q)\big)$. The score is a signed nearest-pattern margin: a contract that resembles known theft code more than known wallet code scores above $\tfrac12$.

The exclusion $F_j\ne F_q$ makes training-time and test-time behaviour identical. In the zero-day protocol a test family is absent from the memory by construction; excluding the query's own family whenever the memory is queried for a training contract---for example when selecting the operating threshold or analysing score distributions---prevents a contract from matching itself or its own clones with similarity $1.0$, a value genuinely new families never attain. With exclusion, the largest similarity a training contract attains in our data is $0.61$, in line with what unseen families exhibit.

The score depends only on the delegate's code. It is therefore available before signing, unaffected by who submits the authorization or how it is batched, and---through canonicalisation---unaffected by the disguises in $\mathcal{M}$. Its cost is $O(|\mathcal{M}|\,|G|)$ set operations per previously unseen contract, with results cached per canonical code.

\subsection{Case-based explanation}
Every score is accompanied by its evidence: the nearest malicious and nearest benign memory entries, their similarities, and the idioms that differ between the query and each match. Because memory entries are deployed contracts, an analyst can inspect the matched contract, its delegators and its drains directly on chain. This form of explanation is faithful by construction---the score is a function of exactly these two similarities---and does not require a post-hoc approximation of a black-box model.

\subsection{Operating point}
The decision threshold is chosen on validation families for a target false-positive rate of 1\%, computed on benign validation delegations only. Thresholds are never tuned on test data. Because the score depends on the delegate's code alone, one threshold serves every decision time: pre-signature, inclusion and monitoring.

\subsection{Complexity and deployment}
Canonicalisation is linear in code length, and the reachability closure visits each instruction at most once. Scoring a previously unseen contract costs one Jaccard computation per memory entry over sets of a few hundred four-grams; with $60$ labelled training contracts this is sub-millisecond, and results are cached per canonical code. A wallet can therefore score a delegation request locally while the user is reviewing the signature prompt, and update the memory as new families are confirmed, without retraining any model.
